\documentclass[10pt]{article}

\usepackage[T1]{fontenc}
\usepackage[utf8]{inputenc}
\usepackage{lmodern}
\usepackage[letterpaper,margin=0.68in]{geometry}
\usepackage{xcolor}
\usepackage{array}
\usepackage{tabularx}
\usepackage{enumitem}
\usepackage{fancyhdr}
\usepackage{hyperref}
\usepackage{microtype}

\definecolor{Ink}{HTML}{18201D}
\definecolor{Muted}{HTML}{52605A}
\definecolor{LaserGreen}{HTML}{16D96B}
\definecolor{DeepGreen}{HTML}{087A3D}
\definecolor{PaleGreen}{HTML}{E9F8EF}
\definecolor{ShieldCyan}{HTML}{2AB9C8}
\definecolor{PaleCyan}{HTML}{E8F7F8}
\definecolor{DroneAmber}{HTML}{D98324}
\definecolor{PaleAmber}{HTML}{FFF3E2}
\definecolor{RuleGray}{HTML}{CBD4D0}
\definecolor{PanelGray}{HTML}{F2F5F3}

\hypersetup{
  colorlinks=true,
  linkcolor=DeepGreen,
  urlcolor=DeepGreen,
  pdftitle={Freelang Doom Engine v6.11.1 User's Guide},
  pdfauthor={The Dosyago Corporation},
  pdfsubject={Quick-start and field guide for Freelang Doom Engine v6.11.1}
}

\setlength{\parindent}{0pt}
\setlength{\parskip}{5pt}
\setlist[itemize]{leftmargin=1.25em,itemsep=2pt,topsep=3pt}
\setlist[enumerate]{leftmargin=1.4em,itemsep=3pt,topsep=3pt}
\renewcommand{\arraystretch}{1.22}
\newcolumntype{K}{>{\raggedright\arraybackslash\bfseries}p{0.25\linewidth}}
\newcolumntype{Y}{>{\raggedright\arraybackslash}X}

\pagestyle{fancy}
\fancyhf{}
\renewcommand{\headrulewidth}{0.4pt}
\renewcommand{\footrulewidth}{0pt}
\fancyhead[L]{\small\color{Muted} FREELANG DOOM ENGINE}
\fancyhead[R]{\small\color{Muted} v6.11.1 USER'S GUIDE}
\fancyfoot[C]{\small\color{Muted} \thepage}

\newcommand{\key}[1]{%
  \begingroup
  \setlength{\fboxsep}{2.5pt}%
  \colorbox{Ink}{\textcolor{white}{\sffamily\bfseries\footnotesize #1}}%
  \endgroup
}

\newcommand{\commandline}[1]{%
  \begingroup
  \setlength{\fboxsep}{6pt}%
  \colorbox{Ink}{\parbox{\dimexpr\linewidth-2\fboxsep}{%
    \textcolor{LaserGreen}{\ttfamily\small #1}}}%
  \endgroup
}

\newenvironment{guidepanel}[2]{%
  \par\medskip
  \begingroup
  \def\GuidePanelAccent{#1}%
  \color{#1}\hrule height 1.4pt
  \vspace{5pt}%
  \color{Ink}%
  \noindent\begin{minipage}{\linewidth}
}{%
  \end{minipage}%
  \vspace{4pt}%
  \color{\GuidePanelAccent}\hrule height 0.5pt
  \endgroup
  \par\medskip
}

\newcommand{\sectionrule}[1]{%
  \vspace{2pt}
  {\color{#1}\rule{\linewidth}{1.5pt}}
}

\begin{document}

\hypersetup{pageanchor=false}
\begin{titlepage}
\thispagestyle{empty}
\color{Ink}

\vspace*{0.35in}
{\sffamily\bfseries\fontsize{15}{18}\selectfont THE DOSYAGO CORPORATION}\par
\vspace{0.12in}
{\color{LaserGreen}\rule{\linewidth}{4pt}}

\vfill
{\sffamily\bfseries\fontsize{35}{39}\selectfont FREELANG DOOM\par}
{\sffamily\bfseries\fontsize{35}{39}\selectfont ENGINE\par}
\vspace{0.16in}
{\sffamily\fontsize{19}{23}\selectfont v6.11.1 USER'S GUIDE}\par

\vspace{0.42in}
{\color{LaserGreen}\rule{0.82\linewidth}{3pt}}\hspace{5pt}%
{\sffamily\bfseries\fontsize{22}{24}\selectfont\color{LaserGreen} >}\par

\vspace{0.5in}
\begin{tabularx}{\linewidth}{>{\sffamily\bfseries\Large\color{DeepGreen}}X >{\sffamily\bfseries\Large\color{DroneAmber}}X >{\sffamily\bfseries\Large\color{ShieldCyan}}X}
LASER & FPV & SHIELD \\
\end{tabularx}

\vfill
\begin{guidepanel}{LaserGreen}{PaleGreen}
\sffamily\bfseries EXPLORER FIELD EDITION\par
\normalfont
A compact guide to launching the engine, surviving the first room, flying the
attack drone, and using the permanent exploration shield. Bring your own
compatible Doom-format WAD.
\end{guidepanel}

{\small\color{Muted}
Independent engine. No WAD is included or downloaded. This guide grants no
rights to third-party game data or trademarks.}
\end{titlepage}
\hypersetup{pageanchor=true}
\setcounter{page}{1}

\section*{Quick start}
\sectionrule{LaserGreen}

\begin{enumerate}
  \item Download the archive for your computer from the v6.11.1 GitHub release.
  \item Extract it. On macOS or Linux, keep \texttt{doom-play} executable.
  \item Supply a legally obtained, compatible Doom-format IWAD or PWAD.
  \item Launch it from a terminal. The map name is optional.
\end{enumerate}

\commandline{macOS/Linux: ./doom-play /path/to/DOOM.WAD E1M1}
\vspace{4pt}
\commandline{Windows: doom-play.exe C:\textbackslash path\textbackslash to\textbackslash DOOM.WAD E1M1}

\begin{guidepanel}{LaserGreen}{PaleGreen}
\textbf{No map argument?} The engine selects the first complete classic map
namespace and opens its map menu. It never searches for or downloads a WAD.
\end{guidepanel}

\subsection*{Core controls}

\begin{tabularx}{\linewidth}{K Y}
\key{Mouse} & Look while the pointer is captured. Click an uncaptured window to recapture it without firing. \\
\key{Left click} / \key{Ctrl} & Fire the current player weapon. Hold to repeat through its cooldown. \\
\key{W} / \key{S} & Move forward / backward. \\
\key{A} / \key{D} & Strafe left / right. \\
\key{Left} / \key{Right} & Turn without the mouse. \\
\key{E} & Use a supported door, lift, switch, or exit line. \\
\key{Space} & Jump. \\
\key{Tab} & Hold the phosphor tactical scan from the active camera. \\
\key{1}-\key{9} & Select an owned weapon. Slot 8 is the laser; slot 9 is the FPV drone. \\
\key{-} & Select the laser directly. Its key echoes the bolt silhouette. \\
\key{Wheel} / \key{[} \key{]} & Cycle through owned weapons. \\
\key{R} & Reset the current map state. \\
\key{Esc} & Release pointer capture; press again to open or close the map menu. \\
\end{tabularx}

\subsection*{Map menu}

Use \key{Up} / \key{Down} to choose any complete classic map namespace found
in the WAD, then press \key{Return}. Player inventory persists when changing
maps. Focus loss clears retained keys and mouse buttons.

\newpage
\section*{Explorer field gear}
\sectionrule{LaserGreen}

The v6.11.1 explorer loadout is intentionally simple: a commodity sidearm, a
permanent defensive field, and a disposable camera with a large payload.

\begin{guidepanel}{LaserGreen}{PaleGreen}
{\sffamily\bfseries\large 8 / LASER BLASTER}\par
\textbf{What it is:} a recoil-free pistol-scale weapon whose exceptional part
is the bright green bolt, not an elaborate body. The pearl-white pistol keeps
the original hand pose and a small green muzzle flare.\par
\textbf{How to use it:} collect the entrance pickup, press \key{8} or \key{-},
aim with the exact camera ray, and fire. It has unlimited charge, a fast rate,
and railgun-level direct damage. Pickups for lesser weapons do not displace it.
\end{guidepanel}

\begin{guidepanel}{ShieldCyan}{PaleCyan}
{\sffamily\bfseries\large PERMANENT EXPLORATION SHIELD}\par
\textbf{What it is:} a small glowing bracelet pickup near the entrance. Once
collected, the shield remains with the player and does not use an armor meter.\par
\textbf{What it stops:} hitscan fire, energy/projectile attacks, Lost Soul
charges, and explosions.\par
\textbf{What still hurts:} real close-range melee attacks and environmental
hazards. The field is permanent in this release; it does not randomly degrade.
\end{guidepanel}

\begin{guidepanel}{DroneAmber}{PaleAmber}
{\sffamily\bfseries\large 9 / FPV ATTACK DRONE}\par
\textbf{What it is:} a single-use flying weapon with a remote camera and a
large 768-damage, 240-radius payload. Every level starts you with ten; extra
pickups are placed from monster-cluster evidence and vary on each load.\par
\textbf{How to launch it:} press \key{9}, then fire once. The camera transfers
to the drone. Your body remains in the world and can still be attacked.
\end{guidepanel}

\subsection*{A useful rhythm}

Take the shield first, select the laser for ordinary exploration, and save the
drone for dense rooms or large monsters. The shield protects your unattended
body from ranged fire while piloting, but a monster that reaches melee range
can still hurt you.

\newpage
\section*{FPV flight school}
\sectionrule{DroneAmber}

The drone carries persistent velocity with light periodic drag. Inputs add
thrust rather than moving a fixed distance, so plan turns early and use
counter-thrust when you need to slow down quickly.

\begin{tabularx}{\linewidth}{K Y}
\key{Mouse} & Aim the wide-angle FPV camera. Forward thrust follows its exact center line. \\
\key{W} / \key{S} & Add forward / reverse thrust. \\
\key{A} / \key{D} & Add lateral thrust. \\
\key{Left} / \key{Right} & Turn; the view banks visibly into the turn. \\
\key{Space} & Climb. \\
\key{Ctrl} & Descend. Control is not fire while the FPV link is active. \\
\key{E} & Use doors and supported switches from the drone's position. \\
\key{Tab} & Hold the phosphor tactical scan from the drone's position and altitude. \\
\key{X} & Arm the impact fuse for this flight. Arming is one-way. \\
\key{Click} / \key{F} & Manually detonate. Release before a fresh click can deploy another drone. \\
\end{tabularx}

\begin{guidepanel}{DroneAmber}{PaleAmber}
\textbf{Contact rules}\par
Every launch begins \textbf{SAFE}. While safe, touching geometry or a monster
stops the drone without detonating it, at any speed. Press \key{X} when you want
the next real geometry or monster contact to trigger the payload. A false sky
ceiling never triggers it or imposes a timer, and the drone can remain airborne
as long as you want. Windows and other legitimate map openings are traversable
when the map permits the drone's collision volume.
\end{guidepanel}

\subsection*{Three reliable approaches}

\begin{itemize}
  \item \textbf{Door scout:} slow before the door, tap \key{E}, then add forward
  thrust only after the opening is clear.
  \item \textbf{Window entry:} center the reticle in the opening, build modest
  speed, and avoid lateral input until the body is through.
  \item \textbf{Payload run:} line up on a cluster, accelerate on the camera ray,
  press \key{X} to arm only when committed, then make contact or use click/F
  when the targets fill the reticle.
\end{itemize}

\subsection*{If the drone seems to disappear}

The FPV view ends after detonation and returns to the player. An unarmed impact
can leave the drone stopped against geometry; counter-thrust or turn away. If
your player body dies while you are flying, the session ends normally.

\newpage
\section*{Options, audio, and troubleshooting}
\sectionrule{ShieldCyan}

\subsection*{Useful launch options}

\begin{tabularx}{\linewidth}{K Y}
\texttt{--resolution=WxH} & Choose a deterministic 8:5 logical raster from 320x200 through 1280x800. \\
\texttt{--no-music} & Skip music loading, synthesis, cache access, and playback. \\
\texttt{--diag} & Write sampled frame, shot, and collision diagnostics to standard error. \\
\texttt{--input-log=PATH} & Save a bounded input, movement, blocker, and FPV lifecycle trace for reproduction. \\
\texttt{--demo} & Run the deterministic E1M1 or E1M2 evidence drive. \\
\end{tabularx}

\subsection*{Audio expectations}

macOS packages include the native speaker sidecar. The first map-music render
can take several seconds; later launches reuse a validated derived WAV cache.
Windows and Linux v6.11.1 continue silently because those releases do not yet
admit a native speaker implementation. Gameplay and visual effects still work.

\begin{guidepanel}{ShieldCyan}{PaleCyan}
\textbf{Signed release}\par
The macOS arm64 and x64 executables and their embedded presenter/speaker
components are Developer ID signed; both release archives were accepted by
Apple's notarization service. The Windows executable and embedded presenter are
Azure Artifact Signing signed with an RFC 3161 timestamp. Linux is a static
x86-64 ELF covered by the release checksum and internal release evidence.
\end{guidepanel}

\subsection*{Fast fixes}

\begin{tabularx}{\linewidth}{K Y}
No window & Confirm the WAD path exists and names a compatible Doom-format file. Launch from a terminal to read the error. \\
No sound & Expected on Windows/Linux v6.11.1. On macOS, try \texttt{--no-music} to separate presentation from synthesis/cache issues. \\
Mouse not turning & Click the game window to capture the pointer. Press Escape once to release it. \\
Movement feels stuck & Release all keys and refocus the window; focus changes clear retained input. \\
Map does not appear & Only complete classic ExMy namespaces found in the supplied WAD are listed. Unsupported mechanics may remain inert. \\
Need a repro & Add \texttt{--input-log=PATH}; the bounded trace records player and drone transitions without containing WAD data. \\
\end{tabularx}

\subsection*{Scope and rights}

Freelang Doom Engine is an independent, incomplete Doom-format implementation,
not a source port and not a claim of complete vanilla behavior. Save games,
multiplayer, automatic campaign progression, and some later-map mechanics are
outside this version. No WAD or third-party game asset is included.

The license distributed with the release controls Freelang-authored material.
This guide does not create a commercial compiler license, grant rights to
third-party assets, or promise a production-support level. DOOM and related
marks belong to their respective owners; the project is not affiliated with or
endorsed by id Software, Bethesda Softworks, or ZeniMax Media.

\vfill
{\color{LaserGreen}\rule{\linewidth}{3pt}}
\begin{center}
{\sffamily\bfseries\large FPV / LASER / SHIELD}\par
{\small\color{Muted} Explore first. Fight when it is interesting.}
\end{center}

\end{document}
